\chapter{Introduction to APIs and Python Requests}

\section{Introduction: Real-Life Hook}

Think about things students do online every day:

\begin{itemize}
    \item Checking weather on a phone app  
    \item Seeing the latest news updates  
    \item Searching for movies or books  
\end{itemize}

All these apps get information from \textbf{APIs (Application Programming Interfaces)}. Python allows us to \textbf{access these APIs easily} using the \texttt{requests} library.

\section{Concept: What is an API?}

\begin{definition}{API}
An API (Application Programming Interface) is a \textbf{way for programs to communicate} with other programs or services over the internet.
\end{definition}

\begin{itemize}
    \item Provides data in a \textbf{structured format (usually JSON)}  
    \item Allows your program to get or send information without knowing how the service works internally  
\end{itemize}

\section{Setting Up Python Requests}

Install the requests library:

\begin{verbatim}
pip install requests
\end{verbatim}

\section{Making a Simple API Request}

Example: Fetch data from a public API:

\begin{verbatim}
import requests

response = requests.get("https://api.coindesk.com/v1/bpi/currentprice.json")
data = response.json()

print("Bitcoin price in USD:", data["bpi"]["USD"]["rate"])
\end{verbatim}

- \texttt{requests.get(url)} sends a GET request  
- \texttt{response.json()} parses JSON data  
- You can access keys in the JSON to extract values  

\section{Handling API Errors}

\begin{verbatim}
response = requests.get("https://api.coindesk.com/v1/bpi/currentprice.json")

if response.status_code == 200:
    data = response.json()
    print("Bitcoin price:", data["bpi"]["USD"]["rate"])
else:
    print("Failed to get data. Status code:", response.status_code)
\end{verbatim}

- \texttt{status\_code} shows if the request succeeded (200 = OK)  
- Always check for errors before processing data  

\section{Mini-Project: Weather App}

Concept:

\begin{itemize}
    \item User enters a city name  
    \item Fetch weather data from an API  
    \item Display temperature, description, and humidity  
\end{itemize}

Example using OpenWeatherMap API:

\begin{verbatim}
import requests

api_key = "YOUR_API_KEY"
city = input("Enter city name: ")

url = f"http://api.openweathermap.org/data/2.5/weather?q={city}&appid={api_key}&units=metric"
response = requests.get(url)

if response.status_code == 200:
    data = response.json()
    temp = data["main"]["temp"]
    desc = data["weather"][0]["description"]
    humidity = data["main"]["humidity"]
    print(f"Weather in {city}: {desc}, Temp: {temp}°C, Humidity: {humidity}%")
else:
    print("City not found or API error")
\end{verbatim}

- Students learn \textbf{user input, API request, JSON parsing, and displaying results}  

\section{Step-by-Step Diagram: API Workflow}

\begin{center}
\begin{tikzpicture}[node distance=1.5cm]
\node (start) [draw, ellipse, fill=green!20] {Start Program};
\node (request) [draw, rectangle, fill=yellow!20, below of=start] {Send HTTP Request to API};
\node (response) [draw, rectangle, fill=orange!20, below of=request] {Receive JSON Response};
\node (parse) [draw, rectangle, fill=yellow!20, below of=response] {Parse JSON Data};
\node (display) [draw, rectangle, fill=orange!20, below of=parse] {Display Data to User};
\node (end) [draw, ellipse, fill=green!20, below of=display] {End Program};
\draw[->] (start) -- (request);
\draw[->] (request) -- (response);
\draw[->] (response) -- (parse);
\draw[->] (parse) -- (display);
\draw[->] (display) -- (end);
\end{tikzpicture}
\end{center}

\section{Key Concepts}

\begin{itemize}
    \item API allows programs to \textbf{communicate and fetch data}  
    \item \texttt{requests} library sends HTTP requests  
    \item JSON is a common data format for APIs  
    \item Always handle errors with \texttt{status\_code}  
    \item Students can create apps like \textbf{weather, crypto price, or news apps}  
\end{itemize}

\section{End-of-Chapter Exercises}

\begin{exercise}{Very Short Questions}
\begin{enumerate}
    \item What is an API?  
    \item Name the Python library used to fetch API data.  
    \item What does \texttt{status\_code} 200 mean?  
\end{enumerate}
\end{exercise}

\begin{exercise}{Short Answer Questions}
\begin{enumerate}
    \item Write Python code to send a GET request to an API and print JSON keys.  
    \item Explain why we need to parse JSON data from APIs.  
    \item Give one real-life example of an API students use daily.  
\end{enumerate}
\end{exercise}

\begin{exercise}{Activity-Based Questions}
\begin{enumerate}
    \item Create a Python app that fetches the latest news headlines from a public news API.  
    \item Modify the weather app to show both Celsius and Fahrenheit temperatures.  
    \item Build a simple cryptocurrency tracker app that shows the current price of Bitcoin and Ethereum.  
\end{enumerate}
\end{exercise}

\section{Teacher Notes}

\marginnote{
\textbf{Learning Objectives:}  
\begin{itemize}
    \item Students understand what APIs are and how they work  
    \item Students can send GET requests and parse JSON responses  
    \item Students can build simple real-world apps using API data  
\end{itemize}

\textbf{Teaching Tips:}  
\begin{itemize}
    \item Start with simple public APIs that don’t require keys  
    \item Show students JSON structure and how to access keys  
    \item Encourage students to experiment with different APIs  
\end{itemize}

\textbf{Assessment:}  
\begin{itemize}
    \item Check correct API request and JSON parsing  
    \item Observe data display and error handling in apps  
\end{itemize}
}